\section{矩阵抽取与傅里叶解码}
\label{sec:matrix}

我们把矩阵测试的高接受率转化为两个策略在稀疏投影问题 $U,O$ 上的
一致性。每个策略只用其自身问题与有界数目的矩阵建议行。所得一致性
概率为正，且与元组长度 $k$ 及奇偶实例无关。表函数只通过其折叠与
共享恒等式以及接受界进入证明，故该证明适用于任何具备这些性质的
族；每个顶点染色的恢复函数都是例子。精确假设在固定参数并定义建议
实验之后陈述。

\subsection{逆短码定理}

结构性输入描述的是在均匀秩一扰动下经常保持取值的函数。这样的函数
在「固定有界数目的行、列线性组合」所得的集合上与一个仿射求值
一致。

记 $\operatorname{Mat}_{\ell,m}=\Hom(\F_2^m,\F_2^\ell)$。
\emph{行与列切片}是由方程
\begin{equation}
 \mathcal S=
 \{M:d_i^\top M=s_i^\top\ (1\le i\le n_r),\quad
          Mq_j=t_j\ (1\le j\le n_c)\}
 \label{eq:shortcode-slice}
\end{equation}
规定的非空集，其中 $d_i\in\F_2^\ell$、$s_i\in\F_2^m$、
$q_j\in\F_2^m$、$t_j\in\F_2^\ell$。计数 $n_r,n_c$ 不必是其方程组
的秩。

\begin{theorem}[二次短码的相等可靠性]
\label{thm:shortcode}
对每个 $\eta\in(0,1)$，存在只依赖 $\eta$ 的 $\alpha\in(0,1]$、
整数 $r_{\rm s}\ge1$ 与整数 $\ell_0$，具有以下性质。对每个
$\ell\ge\ell_0$ 有 $m_0=m_0(\eta,\ell)$，使得对每个 $m\ge m_0$ 与
每个函数 $f:\operatorname{Mat}_{\ell,m}\to\F_2^\ell$，若
\begin{equation}
 \Pr_{M,a,l}\bigl[f(M)=f(M+a l^\top)\bigr]\ge\eta,
 \label{eq:shortcode-acceptance}
\end{equation}
则存在如 Equation~\eqref{eq:shortcode-slice} 的切片 $\mathcal S$
（$n_r,n_c\le r_{\rm s}$）与某个 $z\in\F_2^m$、$u\in\F_2^\ell$，使
\begin{equation}
 \Pr_{M\mid\mathcal S}[f(M)=Mz+u]\ge\alpha.
 \label{eq:shortcode-conclusion}
\end{equation}
此处 $M$、$a\in\F_2^\ell$、$l\in\F_2^m$ 独立均匀，而
$M\mid\mathcal S$ 在切片上均匀。对 $f$ 不施加折叠条件。
\end{theorem}

\begin{proof}
我们从 Khot--Minzer--Safra 的 Grassmann 扩张定理
\cite[Theorem~1.12 in the 2018 revision]{KMS23} 导出此推论，使用
Barak--Kothari--Steurer 的矩阵图卡
\cite[Definition~3.2 and Lemmas~3.3--3.7]{BKS18}。

环境二元空间 $E$ 上的 Grassmann 图以 $\ell$ 维子空间为顶点；两个
顶点相邻当其交的维数为 $\ell-1$。非空顶点集的\emph{保留率}是：
集合中的均匀顶点有一个均匀邻居仍在集合中的概率。所引定理说：保留
率至少固定 $\zeta>0$ 的集合在某个区间
\[
 \{L:A_0\subseteq L\subseteq B_0\},
 \qquad \dim A_0+\operatorname{codim}_E B_0\le r
\]
中有正的相对密度。密度界与 $r$ 只依赖 $\zeta$。定理对所有足够大的
$\ell$、进而所有足够大的 $\dim E$ 成立。

相等接受是各非空纤维 $f^{-1}(y)$ 在矩阵步下保留率的加权平均（权重
为均匀质量）。故某纤维 $S$ 的保留率至少 $\eta$。两个扰动因子之一
为零的概率是
\[
 p_0=2^{-\ell}+2^{-m}-2^{-\ell-m}.
\]
把维数阈值选得使 $p_0\le\eta/2$。在两个因子都非零的条件下，$S$ 的
保留率至少 $(\eta-p_0)/(1-p_0)\ge\eta/2$。

对每个矩阵令
\[
 L_M=\{(x,M^\top x):x\in\F_2^\ell\}
 \subseteq E:=\F_2^\ell\oplus\F_2^m.
\]
映射 $M\mapsto L_M$ 是到图卡 $\mathcal C$（到第一加项投影可逆的
子空间）的双射。由于
\[
 \dim(L_M\cap L_N)=\ell-\rank(M-N),
\]
图卡内的邻居恰对应非零的秩一矩阵差。每个图卡顶点有
$(2^\ell-1)(2^m-1)$ 个这样的邻居，而其完整 Grassmann 度为
$(2^\ell-1)(2^{m+1}-2)$：先在顶点中选超平面，再在该超平面的商中选
一维延拓（非顶点本身）。故恰有一半邻居在图卡内。在 $\F_2$ 上，
每个非零秩一矩阵都有唯一分解 $a l^\top$（因子非零）。因此提升集
$\widetilde S=\{L_M:M\in S\}$ 的 Grassmann 保留率至少 $\eta/4$。

以 $\zeta=\eta/4$ 应用扩张定理，得某个区间上的密度 $\alpha$，其中
$\dim A_0+\operatorname{codim}_E B_0\le r_{\rm s}$。该区间与
$\mathcal C$ 的交只会提高 $\widetilde S$ 的相对密度（因
$\widetilde S\subseteq\mathcal C$），且交非空。$A_0$ 的基
$(d_i,s_i)$ 用方程 $d_i^\top M=s_i^\top$ 刻画
$A_0\subseteq L_M$。对 $E$ 的标准点积，$B_0^\perp$ 的基 $(t_j,q_j)$
用 $Mq_j=t_j$ 刻画 $L_M\subseteq B_0$。二者的联立解集即为所需
形式的切片，两类方程各至多 $r_{\rm s}$ 个。在该切片的至少 $\alpha$
比例上，$f$ 等于定义 $S$ 的那个值。取 $z=0$、$u$ 等于该值即得
结论。

常数 $\alpha,r_{\rm s}$ 的选取只用 $\eta$。由于
$\dim E=\ell+m$，固定 $\ell$ 后剩余的维数阈值被吸收进
$m_0(\eta,\ell)$。未用任何折叠条件。纤维论证是
\cite[Lemma~2.3]{BKS18} 的相等推论；该文的编号引用均指其 arXiv
版本 1。
\end{proof}

Theorem~\ref{thm:shortcode} 中的常数与维数阈值可以有效地选取，
当 $\eta$ 有理时 $\alpha$ 亦有理。这由 \cite[Theorem~2.13 and
Section~2.4]{KMS23} 的定量界与 \cite[Lemmas~2.7--2.8]{KMS23} 的
比较估计得到；定理编号同样指 2018 修订版。

现在固定下文所需参数。从 Lemma~\ref{lem:latent} 选 $p_*>0$ 与相应
的 $r_*$。令
\begin{equation}
 \eta=2^{-r_*}/16,
 \qquad
 2^{-(\ell-r_{\rm s})}<\alpha/8,
 \qquad
 2^{-\ell}<\eta/4,
 \label{eq:matrix-parameters}
\end{equation}
其中 $\alpha,r_{\rm s}$ 先由 Theorem~\ref{thm:shortcode} 给出，
然后选足够大的 $\ell$ 满足该定理、潜在引理以及所示不等式。在此
$\ell$ 值处取潜在数据 $(\mathcal V,K,C,\mu)$，并固定线性分裂
$\mathcal V=K\oplus W$。所有这些对象都在 $k$ 增长之前固定。
特别地，$\ell>r_{\rm s}$。我们把表写成 $P=(M,T)$，其中
$M:L_U\to K$、$T:L_U\to W$。

\subsection{建议实验与解码陈述}

我们先指定两个策略的局部输入，其答案一致性将与外层博弈间隙矛盾。
该实验对每个正整数 $k$ 与每个 $0<\beta<1$ 定义。设 $\mathcal Q_k$
由所有方程问题 $U\in\mathcal E^k$ 以及由它们经
Section~\ref{sec:game} 的坐标投影得到的容许投影问题 $O$ 组成。
问题 $U,O$ 与投影 $\bar\pi:L_U\to L_O$ 的律见 Section~\ref{sec:game}：
在每个块中独立地，方程以概率 $1-\beta$ 被保留，均匀选取的变量位置
以概率 $\beta$ 被保留。

\begin{definition}[投影建议实验]
\label{def:projected-advice}
采样 $U,O,\pi$，再采样独立的均匀映射 $A:K\to\F_2^{r_{\rm s}}$ 与
独立均匀映射
\[
 T_O:L_O\to W,\qquad M_O:L_O\to K.
\]
令 $S_O=AM_O$，并把以下输入发给两个证明者：
\begin{equation}
 \begin{aligned}
  &\text{$U$ 证明者：}&& (U,A,T_O\bar\pi,S_O\bar\pi),\\
  &\text{$O$ 证明者：}&& (O,A,T_O,S_O).
 \end{aligned}
 \label{eq:projected-advice}
\end{equation}
所给建议之外的映射 $M_O$ 被隐藏。每个证明者返回一个实际答案，即
其齐次答案空间中首比特为 1 的向量。当其答案 $a_U,a_O$ 满足
$\bar\pi a_U=a_O$ 时成功。
\end{definition}

$U$ 消息包含映射 $(T_U,S_U)=(T_O\bar\pi,S_O\bar\pi)$，由在 $O$ 上
采样的映射补零得到。我们将先在另一种律下找到有用的建议：$U$ 上的
无限制均匀表 $(M,T)$，取 $S_U=AM$。随后在把投影选择平均之后，把
完整的补零表 $(M_O\bar\pi,T_O\bar\pi)$ 与该无限制表比较。这一比较
是从抽取论证通向实验的桥梁。

\begin{proposition}[矩阵解码]
\label{prop:matrix-decoding}
按 Equation~\eqref{eq:matrix-parameters} 指定的顺序选取潜在与短码
参数，并在 $k$ 之前固定 $\mathcal V=K\oplus W$。存在只依赖这些
固定参数的常数 $\gamma>0$ 与整数 $k_0\ge8$，具有以下性质。

对每个立方数 $k\ge k_0$ 与每个奇偶实例，设对每个 $Q\in\mathcal Q_k$
给定了函数
\[
 F_Q:\Hom(L_Q,\mathcal V)\longrightarrow K.
\]
假设对每个这样的 $Q$、每个 $P\in\Hom(L_Q,\mathcal V)$ 与每个
$c\in K$ 满足折叠，并对每个容许投影 $\bar\pi:L_U\to L_O$ 与每个
$P'\in\Hom(L_O,\mathcal V)$ 满足共享：
\begin{equation}
 \begin{aligned}
 F_Q(P+c e_0^\top)&=F_Q(P)+c,\\
 F_U(P'\bar\pi)&=F_O(P').
 \end{aligned}
 \label{eq:matrix-family-identities}
\end{equation}
若该族以至少 $0.99$ 的概率通过 Equation~\eqref{eq:matrix-test} 的
表测试，则取 $\beta=k^{-2/3}$ 的投影建议实验具有局部策略，其完整
答案一致的概率至少 $\gamma$。
\end{proposition}

$\mathcal G$ 顶点的每个染色都给出这样的族：其恢复函数由
Equations~\eqref{eq:table-folding} 与~\eqref{eq:key-sharing} 满足
Equation~\eqref{eq:matrix-family-identities}。抽象表述记录的是：
解码论证不使用顶点构造的任何其他性质。在定义两个应答规则时该族
已固定，并可以内置于规则之中。

接下来两个小节对任意正整数 $k$ 与满足
Equation~\eqref{eq:matrix-family-identities} 的固定族成立。它们
使用 $U$ 上的无限制均匀表，尚未限制 $\beta$。在把该律与投影实验
比较时才特化 $k$ 与 $\beta$。

\subsection{比较两种矩阵扰动}

均匀矩阵扰动具有 \cite[Lemma~2.10 in the 2018 revision]{KMS23}
计算的按秩索引的傅里叶谱。我们把它与潜在 gadget 提供的噪声比较。

\begin{lemma}[谱比较]
\label{lem:matrix-spectral}
若表族以至少 $0.99$ 的概率通过 Equation~\eqref{eq:matrix-test}，则
\begin{equation}
 \E_{U,T}\Pr_{M,a,l}
   [F_U(M,T)=F_U(M+a l^\top,T)]
 \ge 0.9\,2^{-r_*},
 \label{eq:surrogate-acceptance}
\end{equation}
其中 $a$ 在 $K$ 中均匀，其余矩阵与向量样本在各自空间中均匀。因此
至少 $10\eta$ 的对 $(U,T)$ 的内层概率至少 $4\eta$。
\end{lemma}

\begin{proof}
对 $\sigma\in K^*$ 定义符号函数
$f_{U,\sigma}(P)=(-1)^{\sigma(F_U(P))}$。矩阵域上的频率是泛函元组
$g=(g_i)_{i=0}^{m-1}\in\mathcal V^*$；其在 $P$ 处的特征为
$(-1)^{\sum_i g_i(Pe_i)}$。所有傅里叶系数用均匀概率测度。平移
$a l^\top$ 把该特征乘以 $(-1)^{\sum_i l_i g_i(a)}$。对均匀 $l$ 取
平均得特征值
\[
 \Lambda_g=\Pr_{a\sim\mu}[g_i(a)=0\text{ 对每个 }i].
\]
它在 $[0,1]$ 中。由 Lemma~\ref{lem:latent}，只要
\[
 \dim\Span\{g_i|_K:0\le i<m\}\ge r_*
\]
它就至多为 $7/8$。事实上 $g_i$ 的张成是 $\mathcal V^*$ 的子空间，
引理适用于它。

用一个均匀输出特征展开相等，再用傅里叶正交性，把测试接受率表示为
\[
 \E_{U,\sigma}\sum_g
      \abs{\widehat f_{U,\sigma}(g)}^2\Lambda_g.
\]
Parseval 给出每个符号函数的总质量平方为 1。若 $B$ 是限制秩至少
$r_*$ 上的平均质量平方，则 $0.99\le1-B/8$，故 $B\le0.08$。特别
地，至少 $0.9$ 的平均质量落在低于 $r_*$ 的秩上。

当 $a$ 改为在 $K$ 中均匀时，$g$ 处的特征值为
$2^{-\dim\Span\{g_i|_K\}}$。上述质量因此给出
Equation~\eqref{eq:surrogate-acceptance} 的下界。该扰动固定
$T$，故其中的表达式恰是固定 $(U,T)$ 处普通短码相等测试的平均。
其下界为 $14.4\eta$。设 $q$ 为内层接受至少 $4\eta$ 的概率，则该
平均至多 $4\eta+q$。于是 $q\ge10.4\eta\ge10\eta$。
\end{proof}

\subsection{获得许多有用的行纤维}

逆定理在每个合格 $(U,T)$ 纤维上给出一个结构化切片。$U$ 证明者
关于 $K$ 值矩阵的信息是对 $(A,AM)$。我们需要：这些被观察的行以正
概率决定一个有用的切片。我们还要识别该切片提供的合法答案短集。
下面的擦除论证证明概率界，折叠随后使答案集非空。

对这一抽取步骤，照旧采样 $U$，并取独立均匀映射 $T:L_U\to W$、
$A:K\to\F_2^{r_{\rm s}}$、$M:L_U\to K$。对固定的 $U,T,A$ 与映射
$S_0:L_U\to\F_2^{r_{\rm s}}$，集合 $\{M:AM=S_0\}$ 是其\emph{行
纤维}。称行建议 $(U,T,A,S_0)$ 为\emph{好的}，若存在至多 $r_{\rm s}$
个列规定与某个目标 $Mz+u$，使列规定与 $AM=S_0$ 的联立解集非空，
且在该集合的均匀测度下 $F_U(M,T)=Mz+u$ 以至少 $\alpha/2$ 的概率
成立。

这是行纤维上的解析条件。目标向量 $z$ 起初不必是实际答案。固定的
列值允许 $z$ 在其列方向张成的模下变化，$u$ 相应变化。因此切片
自然提供可能目标向量的一个陪集。其合法部分将是 $U$ 证明者采样的
答案集。

\begin{lemma}[好行建议及其短答案集]
\label{lem:good-row-advice}
对所有足够大的 $k$，矩阵测试接受率至少 $0.99$ 的表族满足
\begin{equation}
 \Pr_{U,T,A,M}[(U,T,A,AM)\text{ 是好的}]
 \ge g_0,
 \qquad
 g_0=10\eta^2\,2^{-\ell r_{\rm s}}>0.
 \label{eq:good-advice-mass}
\end{equation}
该下界与所需的 $k$ 之大与奇偶实例及表族无关。

对每个正整数 $k$ 与每个满足
Equation~\eqref{eq:matrix-family-identities} 的族，可以确定性地为
每个好行建议选取至多 $r_{\rm s}$ 个列规定 $Mq_j=t_j$
（$q_j\in L_U$、$t_j\in K$）与目标 $Mz+u$（$z\in L_U$、$u\in K$、
$e_0^\top z=1$），满足好的定义。记 $Z=\Span\{q_j\}$，所选见证
决定答案集
\begin{equation}
 \mathcal A=(z+Z)\cap\{v\in L_U:e_0^\top v=1\},
 \qquad 1\le |\mathcal A|\le 2^{r_{\rm s}}.
 \label{eq:row-answer-set}
\end{equation}
见证与 $\mathcal A$ 只依赖 $(U,T,A,S_0)$ 与固定的表族。切片及其
目标保持好的定义中的一致性保证。
\end{lemma}

\begin{proof}
固定替代接受率至少 $4\eta$ 的 $(U,T)$，记 $f(M)=F_U(M,T)$。令
$B$ 为矩阵域中所有好行纤维（对一切可能的 $A,S_0$）的并。设其均匀
测度小于 $\eta$。在 $B$ 的表项上修改 $f$，把每一项换成 $K$ 的独立
均匀元素。每个矩阵只被替换一次，即使多条好行纤维在该处重叠。短码
步的两个端点都是均匀矩阵。因此每次这样的修改使接受率至多改变
$2\Pr[B]$，修改后的函数接受率至少 $2\eta$。

我们证明：可以选取随机替换，使其作为与每个目标 $Mz+u$ 匹配的
贡献，在每个行与列的联立切片中少于 $\alpha/4$ 的规模。令
$r=r_{\rm s}$。至多 $r$ 个行规定与 $r$ 个列规定的非空切片，其
齐次空间为
\[
 \Hom\bigl(\F_2^m/\Span\{q_j\},\ \bigcap_i\ker d_i^\top\bigr).
\]
其规模 $n_{\mathcal S}$ 至少为 $2^{(\ell-r)(m-r)}$。用零方程补齐
规定列表后，「切片–目标」描述的个数至多 $2^{(2r+1)(\ell+m)}$。

对固定的描述，只统计随机化表项处的匹配，对其他表项赋以确定的零
贡献。它们是 $[0,1]$ 中的独立变量，总均值至多
$2^{-\ell}n_{\mathcal S}<\alpha n_{\mathcal S}/8$。Hoeffding 不等
式 \cite[Theorem~2, equation~(2.6)]{Hoeffding1963} 把出现至少
$\alpha n_{\mathcal S}/4$ 个这种匹配的概率界定在
$\exp(-\alpha^2 n_{\mathcal S}/32)$ 以内。对所有描述的并界因此
至多
\[
 \exp\left((2r+1)(\ell+m)\log2
       -\frac{\alpha^2}{32}2^{(\ell-r)(m-r)}\right),
\]
当 $m=1+2k$ 增大时趋于零。固定具有所述性质的替换。

考虑任何容许类型的切片。若其行纤维是好的，则该纤维的每个表项都
属于 $B$，其修改后的一致性小于 $\alpha/4$。若行纤维不好，则原
切片上与每个目标的一致性小于 $\alpha/2$。未修改的匹配是原匹配的
子集，而随机化匹配再加少于 $\alpha/4$。修改后的一致性于是小于
$3\alpha/4$。用零补齐较短的行系统表明：这两种情形覆盖
Theorem~\ref{thm:shortcode} 中的每个切片。该定理随即与修改后
函数至少 $2\eta$ 的接受率矛盾。其适用性不要求修改保持折叠。

因此 $\Pr[B]\ge\eta$。对每个 $M\in B$，至少有一个矩阵 $A$ 使某个
好行纤维含 $M$。均匀选取的 $A$ 命中该特定矩阵的概率为
$2^{-\ell r_{\rm s}}$。故在这个固定的 $(U,T)$ 处，好数据的概率
至少 $\eta2^{-\ell r_{\rm s}}$。Lemma~\ref{lem:matrix-spectral}
提供至少 $10\eta$ 个这样的对 $(U,T)$，证明了
Equation~\eqref{eq:good-advice-mass}。所有计数与集中界只依赖固定
参数与 $m$，如所需。

其余断言不使用高测试接受率或 $k$ 之大。固定任意正整数 $k$、满足
Equation~\eqref{eq:matrix-family-identities} 的族与任意好行建议。
固定一个见证，把其列规定写成 $Mq_j=t_j$，令 $Z=\Span\{q_j\}$。
切片非空使规定的值相容：它们确定线性映射 $t:Z\to K$，在整个切片
上 $t(w)=Mw$。把 $z$ 换成 $z+w$、$u$ 换成 $u+t(w)$（$w\in Z$）
因此保持目标。只需证明陪集 $z+Z$ 含首比特为 1 的向量。

否则 $z$ 与 $Z$ 中每个向量首比特都为零。令 $H=\ker A$。对每个
$h\in H$，平移 $M\mapsto M+h e_0^\top$ 保持行纤维、列规定与目标。
由折叠它把 $F_U(M,T)$ 平移 $h$。该作用自由，且其每条 $H$-轨道中
至多一个矩阵满足目标等式。一致性至多
\[
 |H|^{-1}=2^{-\ell+\rank A}
 \le2^{-(\ell-r_{\rm s})}<\alpha/8,
\]
与好性矛盾。选 $w\in Z$ 使 $z+w$ 首比特为 1，并对 $z,u$ 作相应
改变。它们保持原一致性比例。用有限可能性上的固定顺序选一个见证
与这样的代表。该选择只用行建议与固定表族。其向量 $z$ 落在
$\mathcal A$ 中，故集合非空；由 $\dim Z\le r_{\rm s}$，其规模
至多 $2^{r_{\rm s}}$。
\end{proof}

\subsection{转移所选切片}

与稀疏投影问题的比较遵循 Khot--Safra 的光滑外层博弈方法
~\cite[Section~3 and Lemma~3.1]{KS13}，共享线性建议如
Khot--Minzer--Safra~\cite[Sections~3.3--3.4]{KMS17}。
Lemma~\ref{lem:good-row-advice} 给出正概率的行建议：在 $U$ 上的
无限制采样下决定一个切片与一个短的合法答案集。在实际实验中，映射
在 $O$ 上采样并补零到 $U$。我们把完整的补零表（含仿射截距）与
无限制表比较（投影选择被平均掉），随后把所选切片转移到 $O$。

从现在起取大于 1 的立方数 $k$，并在 Definition~\ref{def:projected-advice}
中令 $\beta=k^{-2/3}$。

取 $\beta=k^{-2/3}$ 有两个目的：$k\beta^2\to0$ 使补零表律在
Lemma~\ref{lem:projected-tv} 中逼近无限制律，而 $k\beta\to\infty$
使单例坐标的期望数目发散。一旦建议空间固定，每个单例都有一个一致
为正的概率：每条可见建议行在其变量比特上的系数为零。
Lemma~\ref{lem:clean} 用这些坐标界定一致性。

\begin{lemma}[补零矩阵渐近均匀]
\label{lem:projected-tv}
在每个 $U$ 的条件下，分布
\[
 (M,T)=(M_O\bar\pi,T_O\bar\pi)
\]
与 $\Hom(L_U,K)\times\Hom(L_U,W)$ 上的无限制均匀分布的全变差距离
为 $o(1)$（$k$ 沿立方数趋于无穷）。该界对 $U$ 与奇偶实例一致。
投影选择在此比较中被平均掉。
\end{lemma}

\begin{proof}
把两个分量合成一张 $\mathcal V$ 值表，令 $n=|\mathcal V|$。在未
投影的方程块上，两个斜率是 $\mathcal V$ 中独立的均匀元素。在单例
块上，补零后的斜率对是
\[
 (a,0),\qquad (0,a),\qquad (a,a),
 \qquad a\text{ 在 }\mathcal V\text{ 中均匀}
\]
的等权混合。称这两个对律为 $P_0$ 与 $Q_0$。似然比 $dQ_0/dP_0$ 在
$(0,0)$ 处为 $n$，在三条所示直线上的其余 $3(n-1)$ 个点处为
$n/3$，其他处为零。于是
\[
 \chi^2(Q_0\Vert P_0)
 :=\E_{P_0}\left[\left(\frac{dQ_0}{dP_0}-1\right)^2\right]
 =\frac{n-1}{3}.
\]
因此单块混合 $(1-\beta)P_0+\beta Q_0$ 相对 $P_0$ 的卡方散度为
$\beta^2(n-1)/3$。各块独立采样（即使其问题 ID 重合）。把似然比
相乘给出散度
\begin{equation}
 \left(1+\frac{n-1}{3}\beta^2\right)^k-1=o(1),
 \label{eq:projected-chi-square}
\end{equation}
因为 $n$ 固定而 $k\beta^2=k^{-1/3}$ 趋于零。全变差至多该表达式
平方根的一半。

对方程 $e$ 投影到第三比特的情形，补零截距获得项 $b_e a$（$a$ 为
该单例的斜率）。原截距独立且在 $\mathcal V$ 中均匀。加上这些项
之和后，即使在所有投影选择与斜率都已给出之后，它仍均匀独立。前面
的斜率估计因此证明了整张表的断言。斜率律与
Equation~\eqref{eq:projected-chi-square} 的界都不依赖方程 ID 或其
右端。
\end{proof}

暂时固定好的 $U$-建议及其来自 Lemma~\ref{lem:good-row-advice} 的
所选见证与答案集 $\mathcal A$。对容许投影令 $z'=\bar\pi z$ 与
$Z'=\bar\pi Z$。由于 $\bar\pi$ 保持首坐标，
\begin{equation}
 \bar\pi\mathcal A
 =(z'+Z')\cap\{v\in L_O:e_0^\top v=1\}.
 \label{eq:projected-answer-set}
\end{equation}
事实上，右端陪集中的每个向量在 $z+Z$ 中有原像，且原像有相同的
首坐标。我们将证明 $O$ 证明者的私有傅里叶样本到达这个投影后的
答案集。

见证切片并不自动在投影下存活。例如列条件 $Mq=t$ 变成
$M_O\bar\pi q=t$；若 $\bar\pi q=0$ 而 $t\ne0$ 则不相容。更一般
地，投影后的列方向可能出现新的依赖。下一个论证使用完整表律的
比较，找出正概率的样本：其中所选切片保持非空并保住足够的一致性。

为这一分析，记 $D=(U,O,\pi,A,T_O,S_O)$ 为同时决定两个证明者消息
与投影的联合数据。这些数据在约束 $AM_O=S_O$ 之外不指定 $M_O$。

\begin{lemma}[所选见证的转移]
\label{lem:projected-witness}
设矩阵测试接受率至少 $0.99$。对所有足够大的立方数 $k$，
Definition~\ref{def:projected-advice} 中概率至少为
\begin{equation}
 \gamma_1=\frac{\alpha g_0}{8}\,2^{-\ell r_{\rm s}}>0
 \label{eq:projected-witness-mass}
\end{equation}
的联合数据 $D$ 具有如下性质：补零后的 $U$-建议是好的，带所选
答案集 $\mathcal A$ 及其见证。记 $z'=\bar\pi z$，该见证转移到
非空切片
\[
 AM_O=S_O,\qquad M_O\bar\pi q_j=t_j,
\]
且在其上
\begin{equation}
 \Pr\bigl[F_O(M_O,T_O)=M_Oz'+u\mid
           AM_O=S_O,\ M_O\bar\pi q_j=t_j\text{ 对所有 }j\bigr]
 \ge\alpha/4.
 \label{eq:transferred-witness}
\end{equation}
此式中的概率取遍指定切片上的均匀矩阵。
\end{lemma}

\begin{proof}
先用 $U,A,M,T$ 的无限制采样，令 $S_0=AM$。设 $\mathcal J$ 为事件
「$(U,T,A,S_0)$ 是好的且 $M$ 满足其所选列规定」，$\mathcal M$
为与所选目标相等（在坏数据之外声明该事件为假）。在好行纤维内，
非空列切片施加至多 $\ell r_{\rm s}$ 个额外的标量线性方程。故
Lemma~\ref{lem:good-row-advice} 给出
\[
 \Pr[\mathcal J]\ge g_0\,2^{-\ell r_{\rm s}},
 \qquad
 \Pr[\mathcal J\cap\mathcal M]\ge(\alpha/2)\Pr[\mathcal J].
\]
从而
\begin{equation}
 \Pr[\mathcal J\cap\mathcal M]
       -(\alpha/4)\Pr[\mathcal J]
 \ge\frac{\alpha g_0}{4}\,2^{-\ell r_{\rm s}}.
 \label{eq:witness-positive-difference}
\end{equation}

该表达式中的两个事件只依赖 $(U,A,M,T)$。把独立的 $A$ 加入
Lemma~\ref{lem:projected-tv} 不改变其变差界。因此可以用补零实验
中的值比较两个事件概率。若全变差距离为 $v_k=o(1)$，
Equation~\eqref{eq:witness-positive-difference} 中的差至多改变
$(1+\alpha/4)v_k$。对所有足够大的 $k$，补零后的差至少 $\gamma_1$。

现在对 $D$ 条件化。剩余的 $M_O$ 在行纤维 $AM_O=S_O$ 上均匀。
补零后的 $U$-建议（从而其所选见证与答案集）固定。在好数据上，令
$p_D$ 为投影后列切片的条件概率；当 $p_D>0$ 时，令 $q_D$ 为其
条件相等比例。在坏数据上置 $p_D=0$，当 $p_D=0$ 时置 $q_D=0$。
共享给出 $F_U(M_O\bar\pi,T_O\bar\pi)=F_O(M_O,T_O)$，故补零后的差
恰为
\[
 \E_D\bigl[p_D(q_D-\alpha/4)\bigr]\ge\gamma_1.
\]
被积项在「好性、$p_D>0$、$q_D\ge\alpha/4$」组成的事件之外非正，
且处处至多 1。因此该事件概率至少 $\gamma_1$，且恰具有所述性质。
律比较只用于仅依赖 $(U,A,M,T)$ 的事件；此处对 $D$ 的条件化只是
指出它们的正差在投影实验中的位置。
\end{proof}

\subsection{解码行建议}

从傅里叶质量抽取外层策略的方法遵循 H\aa stad
~\cite[Section~1.1]{Hastad01} 与 Khot--Safra 的傅里叶列表方法
~\cite[Section~1.1]{KS13}。此处投影后的证明者在自己的行纤维上采样
傅里叶质量平方。转移来的切片将表明：该私有样本以正概率到达投影后
的答案集 $\bar\pi\mathcal A$。

现在由各自的消息定义两个策略。联合数据与转移来的见证只用于分析
其一致性；$O$ 证明者不使用 $U$ 证明者的见证或答案集。

$U$ 证明者用其可见的行建议执行 Lemma~\ref{lem:good-row-advice}
中的确定性选择。当建议是好的时，它从 $\mathcal A$ 中均匀采样
$a_U$。这是至多 $2^{r_{\rm s}}$ 个实际答案的非空集。在其他输入上
返回任意实际答案。

$O$ 证明者从自己的行纤维出发。给定 $(O,A,T_O,S_O)$，令 $H=\ker A$，
并按固定规则选满足 $AM_0=S_O$ 的代表 $M_0$。这样的代表在实验的
每个输入上都存在。把行纤维写成 $M_0+X$，其中 $X=\Hom(L_O,H)$。
采样均匀的 $\sigma\in K^*$，并作符号函数
\[
 f_\sigma(N)=(-1)^{\sigma(F_O(M_0+N,T_O))},\qquad N\in X.
\]
以概率 $|\widehat f_\sigma(\Phi)|^2$ 采样傅里叶频率
$\Phi\in X^*$。把 $X^*=H^*\otimes L_O$ 等同，纯张量上的配对定义为
\[
 \langle\varphi\otimes v,N\rangle=\varphi(Nv).
\]
对 $h\in H$，收缩把 $\Phi$ 送到向量 $\Phi(h)\in L_O$。
若 $\sigma|_H\ne0$，按固定规则选使 $\sigma(h)=1$ 的 $h\in H$ 并
返回 $\Phi(h)$。若限制平凡，返回任意实际答案。

该规则在正质量频率处总返回实际答案。事实上
Equation~\eqref{eq:matrix-family-identities} 的折叠恒等式给出
\[
 f_\sigma(N+h e_0^\top)=-f_\sigma(N)
 \qquad\text{当 }\sigma(h)=1.
\]
在傅里叶系数中换变量可知其非零要求
$(-1)^{\langle\Phi,h e_0^\top\rangle}=-1$，等价于
$e_0^\top\Phi(h)=1$。

Figure~\ref{fig:advice} 概括两个解码器的信息与私有选择。转移来的
切片只用于分析其一致概率。
\begin{figure}[tbp]
\centering
\begin{tikzpicture}[
  x=1cm,y=1cm,
  every node/.style={font=\small},
  box/.style={draw=linkblue!80!black,rounded corners=2pt,
    text width=6.0cm,align=center,inner sep=8pt},
  input/.style={box,minimum height=16mm},
  rule/.style={box,minimum height=35mm},
  output/.style={box,minimum height=10mm},
  flow/.style={-{Latex[length=2.1mm]},thick,draw=black!70}
]
\node[box,text width=13.3cm,minimum height=12mm] (sample) at (3.7,0)
  {\textbf{裁判}\\[2pt]
   采样 $U,O,\bar\pi,A,T_O,M_O$；令 $S_O=AM_O$。
   只发送下述输入。};

\node[input] (uinput) at (0,-1.9)
  {\textbf{方程解码器}\\
   $(U,A,T_U,S_U)$\\
   $T_U=T_O\bar\pi,\quad S_U=S_O\bar\pi$};
\node[input] (oinput) at (7.4,-1.9)
  {\textbf{投影解码器}\\
   $(O,A,T_O,S_O)$};
\draw[flow] (sample.south -| uinput.north) -- (uinput.north);
\draw[flow] (sample.south -| oinput.north) -- (oinput.north);

\node[rule] (urule) at (0,-4.8)
  {从本解码器的可见建议中选取固定的切片见证\\及其合法答案集\\[2pt]
   $\mathcal A=(z+Z)\cap\{v:e_0^\top v=1\}$。\\[6pt]
   \textit{私有随机性：}\\
   从 $\mathcal A$ 中均匀采样 $a_U$。};
\node[rule] (orule) at (7.4,-4.8)
  {令 $H=\ker A$。选 $M_0$\\使 $AM_0=S_O$。\\[5pt]
   \textit{私有随机性：}\\
   采样均匀的 $\sigma\in K^*$，再以概率\\
   $|\widehat f_\sigma(\Phi)|^2$ 采样 $\Phi$。\\[5pt]
   若 $\sigma|_H\ne0$，取固定的 $h\in H$\\
   使 $\sigma(h)=1$；令 $a_O=\Phi(h)$。};
\draw[flow] (uinput) -- (urule);
\draw[flow] (oinput) -- (orule);

\node[output] (uoutput) at (0,-7.9)
  {$a_U\in L_U,\qquad e_0^\top a_U=1$};
\node[output] (ooutput) at (7.4,-7.9)
  {$a_O\in L_O,\qquad e_0^\top a_O=1$};
\draw[flow] (urule) -- (uoutput);
\draw[flow] (orule) -- (ooutput);

\node[box,text width=8.2cm,minimum height=11mm] (verify) at (3.7,-9.6)
  {\textbf{裁判检验}\quad $\bar\pi a_U=a_O$};
\draw[flow] (uoutput.south) -- (verify.north -| uoutput.south);
\draw[flow] (ooutput.south) -- (verify.north -| ooutput.south);
\end{tikzpicture}
\caption{两个解码器使用独立的随机带。二者都不接收采样到的 $M_O$
或投影 $\bar\pi$；右侧解码器从可见建议中自行选取代表 $M_0$。每个
所示的建议映射都经由齐次坐标包含其仿射截距。转移来的切片只用于
分析这些私有规则：右侧解码器不使用左侧的见证或答案集。若左侧建议
不好，或右侧 $\sigma|_H=0$，该解码器返回任意合法答案。}
\label{fig:advice}
\end{figure}

\FloatBarrier

\begin{lemma}[转移切片给出一致]
\label{lem:fourier-coset-decoding}
固定满足 Lemma~\ref{lem:projected-witness} 的联合数据。在这些数据
的条件下，刚定义的两个规则使用独立随机性，在 $\bar\pi$ 之下一
致的概率至少
\begin{equation}
 2^{-\ell r_{\rm s}}(\alpha/8)^2\,2^{-r_{\rm s}}.
 \label{eq:conditional-decoding}
\end{equation}
\end{lemma}

\begin{proof}
令 $Z'=\bar\pi Z$，故 $\dim Z'\le r_{\rm s}$。行纤维中投影后的列
规定把限制 $N|_{Z'}$ 固定为线性映射 $b:Z'\to H$。所得仿射切片
非空；取其中一个 $N_*$。对每个 $\sigma\in K^*$ 定义
\[
 C_\sigma=\E\bigl[f_\sigma(N)(-1)^{\sigma(Nz')}
                       \mid N|_{Z'}=b\bigr].
\]
Equation~\eqref{eq:transferred-witness} 的特征展开给出
\[
 \E_\sigma\bigl[(-1)^{\sigma(M_0z'+u)}C_\sigma\bigr]
 \ge\alpha/4,
 \qquad\text{故}\qquad
 \E_\sigma|C_\sigma|\ge\alpha/4.
\]
恰有 $2^{-\dim H}$ 比例的特征在 $H$ 上限制平凡。由
$\dim H\ge\ell-r_{\rm s}$ 与 $|C_\sigma|\le1$，
Equation~\eqref{eq:matrix-parameters} 蕴含
\begin{equation}
 \E_\sigma\bigl[\mathbf 1_{\{\sigma|_H\ne0\}}|C_\sigma|\bigr]
 \ge\alpha/8.
 \label{eq:nontrivial-correlation}
\end{equation}

在 $Z'$ 上为零的矩阵在 $X^*$ 中的零化子是 $H^*\otimes Z'$。令
$\theta_\sigma=(\sigma|_H)\otimes z'$，在仿射切片上的傅里叶展开
给出
\[
 C_\sigma
 =\sum_{\Phi\in\theta_\sigma+H^*\otimes Z'}
       \widehat f_\sigma(\Phi)
       (-1)^{\langle\Phi+\theta_\sigma,N_*\rangle}.
\]
该和有 $2^{\dim H\dim Z'}\le2^{\ell r_{\rm s}}$ 项。
Cauchy--Schwarz（先对每个和、再对 $\sigma$ 的平均）由
Equation~\eqref{eq:nontrivial-correlation} 给出
\begin{align*}
 &\E_\sigma\left[
   \mathbf 1_{\{\sigma|_H\ne0\}}
   \sum_{\Phi\in\theta_\sigma+H^*\otimes Z'}
             |\widehat f_\sigma(\Phi)|^2\right]\\
 &\hspace{15mm}\ge
 2^{-\ell r_{\rm s}}
 \E_\sigma\bigl[\mathbf 1_{\{\sigma|_H\ne0\}}C_\sigma^2\bigr]
 \ge 2^{-\ell r_{\rm s}}(\alpha/8)^2.
\end{align*}
这就是 $O$ 解码器采到非平凡限制与所示陪集中频率的概率下界。

每个这样的频率形如 $\Phi=(\sigma|_H)\otimes z'+\Psi$
（$\Psi\in H^*\otimes Z'$）。对解码器选取的 $h$，
\[
 \Phi(h)=z'+\Psi(h)\in z'+Z'.
\]
如已证，正质量频率给出首比特 1。Equation~\eqref{eq:projected-answer-set}
因此把返回向量置于 $\bar\pi\mathcal A$。独立地，$U$ 解码器从
$\mathcal A$ 均匀采样，以至少 $1/|\mathcal A|\ge2^{-r_{\rm s}}$
的概率选中原像。各界相乘即证
Equation~\eqref{eq:conditional-decoding}。
\end{proof}

\begin{proof}[Proposition~\ref{prop:matrix-decoding} 的证明]
使用上述两个解码器。Lemma~\ref{lem:projected-witness} 提供概率
至少 $\gamma_1$ 的联合数据集，Lemma~\ref{lem:fourier-coset-decoding}
在该集上给出所述条件一致性。于是可取
\[
 \gamma=\gamma_1\,2^{-\ell r_{\rm s}}(\alpha/8)^2\,2^{-r_{\rm s}}>0.
\]
对 $k$ 之大的需要只来自 Theorem~\ref{thm:shortcode}、
Lemma~\ref{lem:good-row-advice} 的集中界与 Lemma~\ref{lem:projected-tv}
的全变差估计。它们对奇偶实例与表族一致。见证选取与傅里叶采样不必
高效：它们在可靠性论证的有限博弈中定义策略。最后，对证明者独立
随机带的取平均也给出一致概率不低于此的确定性策略。
\end{proof}
